% Generated draft: requires mathematical and source review.
\documentclass[11pt]{article}
\usepackage[margin=25mm]{geometry}
\usepackage{fontspec}
\setmainfont{Latin Modern Roman}
\usepackage{amsmath,amssymb,mathtools,unicode-math}
\setmathfont{Latin Modern Math}
\usepackage{enumitem,needspace,xurl,hyperref}
\setlength{\parindent}{0pt}
\setlength{\parskip}{5pt}
\setlength{\emergencystretch}{3em}
\newcommand{\N}{\mathbb{N}}
\newcommand{\Z}{\mathbb{Z}}
\newcommand{\Q}{\mathbb{Q}}
\newcommand{\R}{\mathbb{R}}
\newcommand{\C}{\mathbb{C}}
\newcommand{\F}{\mathbb{F}}
\newcommand{\A}{\mathbb{A}}
\newcommand{\PP}{\mathbb{P}}
\newcommand{\E}{\mathbb{E}}
\newcommand{\eps}{\varepsilon}
\newcommand{\dd}{\,\mathrm{d}}
\begin{document}
\section*{UnsolvedMath ID: 2100702}
\textbf{Source problem code:} catalogue entry
\subsection{Definitions and mathematical statement}
Let $P_1(h),\ldots,P_n(h)$ be the intended semiclassical operator families, with semiclassical parameter $h\downarrow0$ and principal symbols $p_1,\ldots,p_n$ on their classical phase space $M$. The source assumes the preceding conditions labelled \texttt{item:self-adjoint} and \texttt{item:commute}, but their full text is not supplied. They must not be replaced silently by a particular self-adjoint-domain convention or by exact, strong, formal, or asymptotic commutation.

Functional independence of the principal symbols at $x\in M$ means
\[
\operatorname{rank}d(p_1,\ldots,p_n)_x=n,
\quad\text{equivalently}\quad
(dp_1\wedge\cdots\wedge dp_n)(x)\ne0.
\]
The requested independence axiom is that this holds almost everywhere, with respect to the intended classical measure (normally Liouville measure in a symplectic phase space). It is not merely independence somewhere or on one regular fiber.

S. V\~u Ng\d{o}c asks whether this almost-everywhere condition can be detected purely from asymptotics of the joint eigenvalues or eigenfunctions of $(P_1,\ldots,P_n)$. When common eigenvectors exist, their joint point spectrum is
\[
\Sigma_h=\{(\lambda_1,\ldots,\lambda_n):
\exists u\ne0\text{ with }P_j(h)u=\lambda_j u\text{ for every }j\}.
\]
The problem seeks a spectral/asymptotic criterion for the displayed classical rank condition, within the operator class fixed by the missing axioms. No joint discreteness, multiplicity convention, energy window, precision in $h$, or allowed normalization/identification of eigenfunctions is stated. A joint-eigenvalue-only criterion and a criterion using eigenfunctions are different data questions; both alternatives in the source are retained.

The supplied research summary reports the detection question open/unverified. It mentions relations of joint spectral clusters with action variables and quantum monodromy, including work of V\~u Ng\d{o}c and Charles--V\~u Ng\d{o}c, but does not identify a verified theorem answering this precise functional-independence question.
\subsection{Short English statement}
The question asks whether joint semiclassical eigenvalue or eigenfunction asymptotics reveal that the principal symbols have full differential rank almost everywhere. The source's self-adjointness and commutation axioms are referenced but absent, so the operator class and exact spectral data must be recovered before a precise criterion can be posed.
\subsection{Sources}
\begin{itemize}
\item Local UnsolvedMath registry, permanent ID 2100702, exact source code catalogue entry, ``Open Problems in Integrable Systems --- Integrability and Quantisation''; status open, classification OPEN-TRIAGE.
\item \emph{Open Problems, Questions, and Challenges in Finite-Dimensional Integrable Systems} (2018), Problem 7.2, attributed to S. V\~u Ng\d{o}c, \url{https://arxiv.org/abs/1804.03737}. The record identifies extraction from public arXiv source TeX, accessed July 29, 2026, and says the source presented the item as open. The extraction retains the labels \texttt{item:self-adjoint}, \texttt{item:commute}, and \texttt{item:independent} without their definitions.
\item The supplied research summary mentions V\~u Ng\d{o}c's quantum-monodromy work, Charles--V\~u Ng\d{o}c, and semiclassical inverse spectral theory as context. It supplies no titles, theorem numbers, or precise reference resolving the detection problem and expressly reports that it could not verify an answer.
\end{itemize}
\end{document}
